% GENERATED FILE. Edit canonical lessons, then run npm run generate:books.
% canonical-source-hash: fnv1a64:52725295792b90ed
% canonical-lessons: HI-C211-catni, HI-C211-masala, HI-C211-haldi, HI-C211-dhaniya, HI-C211-nibu

\chapter{Chutney, Spice, Turmeric, Coriander, A lemon}
\label{ch:hi-chutney-spice-turmeric-coriander-a-lemon}
\hlchaptermodality{211}

\begin{chapteropening}
\textbf{By the end of this chapter:} \emph{I can say chutney, spice, turmeric, coriander and a lemon.}

Complete the last lesson of chapter 211: i can say chutney, spice, turmeric, coriander and a lemon.
\end{chapteropening}

\section[\texorpdfstring{\hi{चटनी} (caṭnī)}{चटनी (caṭnī)}]{\hi{चटनी} (caṭnī) — chutney}

\label{lesson:HI-C211-catni}



\begin{quote}
Before the new one: say the Hindi for a laddu, then the Hindi for pickle.
\end{quote}

\subsection*{You'll want to know: \hi{चटनी}}
\textbf{\hi{चटनी}} — \emph{caṭnī} — \textquotedblleft{}chutney\textquotedblright{}.

\begin{grammarlens}[title={What you've built}]
One more word for this chapter.
\end{grammarlens}

\subsection*{Your turn}
\begin{itemize}
\raggedright
  \item \emph{Say it:} \emph{caṭnī}
  \item \emph{Say it:} \emph{caṭnī}, once more
  \item \emph{Say it:} say \emph{caṭnī}
  \item \emph{Recall:} say the Hindi for a laddu, then the Hindi for pickle, then say \emph{caṭnī} again
\end{itemize}

\begin{tcolorbox}[breakable,colback=teal!4,colframe=teal!35!black,title={Before you move on}]
What does \hi{चटनी} mean? (\textquotedblleft{}Chutney\textquotedblright{}.) Say it once more.
\end{tcolorbox}
\section[\texorpdfstring{\hi{मसाला} (masālā)}{मसाला (masālā)}]{\hi{मसाला} (masālā) — spice, a spice mix}

\label{lesson:HI-C211-masala}



\begin{quote}
Before the new one: say the Hindi for pickle, then the Hindi for chutney.
\end{quote}

\subsection*{You'll want to know: \hi{मसाला}}
\textbf{\hi{मसाला}} — \emph{masālā} — \textquotedblleft{}spice, a spice mix\textquotedblright{}.

\begin{grammarlens}[title={What you've built}]
One more word for this chapter.
\end{grammarlens}

\subsection*{Your turn}
\begin{itemize}
\raggedright
  \item \emph{Say it:} \emph{masālā}
  \item \emph{Say it:} \emph{masālā}, once more
  \item \emph{Say it:} say \emph{masālā}
  \item \emph{Recall:} say the Hindi for pickle, then the Hindi for chutney, then say \emph{masālā} again
\end{itemize}

\begin{tcolorbox}[breakable,colback=teal!4,colframe=teal!35!black,title={Before you move on}]
What does \hi{मसाला} mean? (\textquotedblleft{}Spice, a spice mix\textquotedblright{}.) Say it once more.
\end{tcolorbox}
\section[\texorpdfstring{\hi{हल्दी} (haldī)}{हल्दी (haldī)}]{\hi{हल्दी} (haldī) — turmeric}

\label{lesson:HI-C211-haldi}



\begin{quote}
Before the new one: say the Hindi for chutney, then the Hindi for spice.
\end{quote}

\subsection*{You'll want to know: \hi{हल्दी}}
\textbf{\hi{हल्दी}} — \emph{haldī} — \textquotedblleft{}turmeric\textquotedblright{}.

\begin{grammarlens}[title={What you've built}]
One more word for this chapter.
\end{grammarlens}

\subsection*{Your turn}
\begin{itemize}
\raggedright
  \item \emph{Say it:} \emph{haldī}
  \item \emph{Say it:} \emph{haldī}, once more
  \item \emph{Say it:} say \emph{haldī}
  \item \emph{Recall:} say the Hindi for chutney, then the Hindi for spice, then say \emph{haldī} again
\end{itemize}

\begin{tcolorbox}[breakable,colback=teal!4,colframe=teal!35!black,title={Before you move on}]
What does \hi{हल्दी} mean? (\textquotedblleft{}Turmeric\textquotedblright{}.) Say it once more.
\end{tcolorbox}
\section[\texorpdfstring{\hi{धनिया} (dhaniyā)}{धनिया (dhaniyā)}]{\hi{धनिया} (dhaniyā) — coriander}

\label{lesson:HI-C211-dhaniya}



\begin{quote}
Before the new one: say the Hindi for spice, then the Hindi for turmeric.
\end{quote}

\subsection*{You'll want to know: \hi{धनिया}}
\textbf{\hi{धनिया}} — \emph{dhaniyā} — \textquotedblleft{}coriander\textquotedblright{}.

\begin{grammarlens}[title={What you've built}]
One more word for this chapter.
\end{grammarlens}

\subsection*{Your turn}
\begin{itemize}
\raggedright
  \item \emph{Say it:} \emph{dhaniyā}
  \item \emph{Say it:} \emph{dhaniyā}, once more
  \item \emph{Say it:} say \emph{dhaniyā}
  \item \emph{Recall:} say the Hindi for spice, then the Hindi for turmeric, then say \emph{dhaniyā} again
\end{itemize}

\begin{tcolorbox}[breakable,colback=teal!4,colframe=teal!35!black,title={Before you move on}]
What does \hi{धनिया} mean? (\textquotedblleft{}Coriander\textquotedblright{}.) Say it once more.
\end{tcolorbox}
\section[\texorpdfstring{\hi{नींबू} (nī̃bū)}{नींबू (nī̃bū)}]{\hi{नींबू} (nī̃bū) — a lemon, a lime}

\label{lesson:HI-C211-nibu}



\begin{quote}
Before the new one: say the Hindi for turmeric, then the Hindi for coriander.
\end{quote}

\subsection*{You'll want to know: \hi{नींबू}}
\textbf{\hi{नींबू}} — \emph{nī̃bū} — \textquotedblleft{}a lemon, a lime\textquotedblright{}.

\begin{grammarlens}[title={What you've built}]
One more word for this chapter.
\end{grammarlens}

\subsection*{Your turn}
\begin{itemize}
\raggedright
  \item \emph{Say it:} \emph{nī̃bū}
  \item \emph{Say it:} \emph{nī̃bū}, once more
  \item \emph{Say it:} say \emph{nī̃bū}
  \item \emph{Recall:} say the Hindi for turmeric, then the Hindi for coriander, then say \emph{nī̃bū} again
\end{itemize}

\begin{tcolorbox}[breakable,colback=teal!4,colframe=teal!35!black,title={Before you move on}]
What does \hi{नींबू} mean? (\textquotedblleft{}A lemon, a lime\textquotedblright{}.) Say it once more.
\end{tcolorbox}
